\subsection{加权最小二乘（WLS）}
\label{sec:regression-wls}

OLS默认各条记录的误差方差相同。实际数据常常不是这样：低收入人群的消费能力有限、消费波动小，高收入人群的消费方差却大；同一台仪器在高量程档的测量噪声也比低量程档大。把每条记录的误差记作$\varepsilon_i$，若$\operatorname{Var}(\varepsilon_i\mid\bm{x}_i)=\sigma_i^2$随$i$变化，就称为异方差。

此时OLS估计仍是无偏的——前提是线性均值设定正确、误差零条件均值——但不再是有效的：精度高的记录（$\sigma_i^2$小）与精度差的记录在目标里地位相同，后者把估计来回拉动，估计方差不是最小的。解决办法是给每条残差平方除以它自己的误差方差——由高斯--马尔可夫定理（加权形式），当各记录误差独立时，使估计方差最小的权重恰与误差方差成反比。这样$n$条记录各自按精度进入目标，预测形式仍保持$f(\bm{x})=b+\bm{w}^{\top}\bm{x}$。

\subsubsection{模型结构}
\label{sec:regression-wls-structure}

独立异方差、且各记录误差方差$\sigma_i^2>0$已知时，\term{加权最小二乘}（weighted least squares, WLS）的目标为
\begin{equation}
  (\widehat{\bm{w}},\widehat b)
  \in\argmin_{\bm{w},b}\ J_{\mathrm{WLS}}(\bm{w},b),
  \qquad
  J_{\mathrm{WLS}}(\bm{w},b)=\frac12\sum_{i=1}^{n}\frac{\bigl(y_i-b-\bm{w}^{\top}\bm{x}_i\bigr)^2}{\sigma_i^2}.
  \label{eq:regression-wls-objective}
\end{equation}
沿用$1/2$约定、不带$1/n$因子，梯度与更新与该约定一致。误差方差大的记录，其平方残差在目标中除以更大的数，训练时对参数的拉动变小；$\sigma_i^2$全部相等时，权重都相同，式\eqref{eq:regression-wls-objective}退回OLS目标。

\subsubsection{参数学习}
\label{sec:regression-wls-learning}

\paragraph{显式梯度与参数更新}

对式\eqref{eq:regression-wls-objective}求导：对$\bm{w}$、$b$分别有
\[
  \nabla_{\bm w}J_{\mathrm{WLS}}
  =\sum_{i=1}^{n}\frac{b+\bm{w}^{\top}\bm{x}_i-y_i}{\sigma_i^2}\,\bm{x}_i,
  \qquad
  \frac{\partial J_{\mathrm{WLS}}}{\partial b}
  =\sum_{i=1}^{n}\frac{b+\bm{w}^{\top}\bm{x}_i-y_i}{\sigma_i^2}.
\]
沿负方向按步长$\eta$迭代：
\begin{align}
  \bm{w}^{(t+1)}
  &=\bm{w}^{(t)}-\eta\sum_{i=1}^{n}\frac{b^{(t)}+\bm{w}^{(t)\top}\bm{x}_i-y_i}{\sigma_i^2}\,\bm{x}_i,\\
  b^{(t+1)}
  &=b^{(t)}-\eta\sum_{i=1}^{n}\frac{b^{(t)}+\bm{w}^{(t)\top}\bm{x}_i-y_i}{\sigma_i^2}.
  \label{eq:regression-wls-gradient}
\end{align}
两式中的$b$与$\bm{w}$都取第$t$轮的值：先按同一轮参数算出加权残差，再同时更新。梯度置零即加权正规方程，可一次解出闭式参数；新输入仍在原始特征上用已学参数计算$f(\bm{x})=b+\bm{w}^{\top}\bm{x}$。
